```latex
\documentclass[11pt]{article}

\usepackage{amsmath,amssymb,amsthm}
\usepackage[margin=1in]{geometry}

\newcommand{\Sph}{\mathbb S^2}
\newcommand{\Chat}{\widehat{\mathbb C}}
\newcommand{\Teich}{\mathcal T}
\newcommand{\D}{\mathcal D}
\newcommand{\Mod}{\operatorname{Mod}}
\newcommand{\PMod}{\operatorname{PMod}}
\newcommand{\id}{\operatorname{id}}

\newtheorem{theorem}{Theorem}[section]
\newtheorem{lemma}[theorem]{Lemma}
\newtheorem{proposition}[theorem]{Proposition}
\theoremstyle{remark}
\newtheorem{remark}[theorem]{Remark}

\title{A Dehn twist acting with an infinite orbit on a deformation space}
\author{}
\date{}

\begin{document}
\maketitle

\section{Statement and result}

Write
\[
p_0=c_1,\qquad p_1=f(c_1),\qquad p_2=f^{\circ 2}(c_1),\qquad
p_3=f^{\circ 3}(c_1).
\]
Thus \(f(p_3)=p_0\), the four points \(p_0,p_1,p_2,p_3\) are
distinct, and
\[
V(f)=\{p_1,v\},\qquad v=f(c_2).
\]
Set
\[
A=\{p_0,p_1,p_2,p_3\},
\qquad
B=A\cup\{v\}.
\]

Choose an embedded arc
\[
\alpha\subset \Sph\setminus\{p_0,p_1,p_3\}
\]
joining \(v\) to \(p_2\), and let \(\Delta\) be a closed regular
neighborhood of \(\alpha\) whose boundary is disjoint from \(B\).
Thus
\[
\Delta\cap B=\{v,p_2\}.
\]
Put
\[
\gamma=\partial\Delta.
\]
Equivalently, \(\gamma\) is a simple closed curve in
\(\Sph\setminus B\) separating
\[
\{v,p_2\}
\quad\text{from}\quad
\{p_0,p_1,p_3\}.
\]
Both complementary components contain at least two points of \(B\), so
\(\gamma\) is essential and nonperipheral.

Let \(T_\gamma\) denote the positive Dehn twist about \(\gamma\).

\begin{theorem}\label{thm:main}
The mapping class \(T_\gamma^2\) preserves \(\D_f\).  Moreover,
for every connected component \(C\) of \(\D_f\), the components
\[
T_\gamma^{2n}(C),\qquad n\in\mathbb Z,
\]
are pairwise distinct.  Hence the action of \(T_\gamma^2\) on
\(\pi_0(\D_f)\) has infinite orbits.
\end{theorem}

The remainder of this note proves Theorem \ref{thm:main}.

\section{The inverse image of \(\gamma\)}

We use the following elementary proper-map version of the
Riemann--Hurwitz formula: if \(F:X\to Y\) is a proper degree-\(d\)
branched cover of compact surfaces, then
\[
\chi(X)=d\,\chi(Y)-\sum_{x\in X}(e_x-1),
\]
where \(e_x\) is the local degree.

\begin{lemma}\label{lem:preimage-disk}
The full preimage \(\widetilde\Delta:=f^{-1}(\Delta)\) is a closed
disk.  Its boundary
\[
\widetilde\gamma:=f^{-1}(\gamma)
\]
is connected and maps to \(\gamma\) with degree \(2\).
\end{lemma}

\begin{proof}
The disk \(\Delta\) contains exactly one critical value, namely
\(v=f(c_2)\); the other critical value is \(p_1\notin\Delta\).
Thus the only ramification point in \(f^{-1}(\Delta)\) is \(c_2\),
with local degree \(2\).

We first show that \(f^{-1}(\Delta)\) is connected.  Every
component of its boundary is a covering of \(\gamma\), and the sum
of the degrees of these coverings is \(2\), the degree of \(f\).
If \(f^{-1}(\Delta)\) were disconnected, it would therefore have
exactly two boundary components, distributed between two components,
each mapping to \(\gamma\) with degree \(1\).  Each such component
would be a disk mapped properly of degree \(1\) onto \(\Delta\),
hence would contain no ramification point.  This contradicts the
fact that \(c_2\in f^{-1}(\Delta)\).  Hence \(f^{-1}(\Delta)\) is
connected.

Riemann--Hurwitz now gives
\[
\chi(f^{-1}(\Delta))
   =2\chi(\Delta)-(2-1)
   =2-1=1.
\]
A connected planar compact surface of Euler characteristic \(1\) is
a closed disk.  Consequently its boundary is connected.  Since the
sum of the boundary mapping degrees is \(2\), the unique component
\(\widetilde\gamma\) covers \(\gamma\) with degree \(2\).
\end{proof}

\begin{lemma}\label{lem:lift-peripheral}
The disk \(\widetilde\Delta=f^{-1}(\Delta)\) contains exactly one
point of \(B\), namely \(p_1=f(c_1)\).  In particular,
\(\widetilde\gamma\) is peripheral in \((\Sph,B)\).
\end{lemma}

\begin{proof}
Since \(v\) is a critical value of multiplicity \(2\) at \(c_2\),
\[
f^{-1}(v)=\{c_2\}.
\]
Moreover \(c_2\notin A\): indeed, if \(c_2=p_j\) for some \(j\), then
\(v=f(c_2)=p_{j+1}\), with indices modulo \(4\), contrary to the
assumption that \(v\) is not in the \(c_1\)-cycle.  Thus \(c_2\)
is not a marked point.

Because \(p_2\) is neither \(p_1\) nor \(v\), it is a regular value.
Hence \(f^{-1}(p_2)\) consists of two distinct points.  One of them
is \(p_1\), since \(f(p_1)=p_2\); the other is unmarked.  Therefore
the only marked point in \(f^{-1}(\Delta)\) is \(p_1\).

It follows that \(\widetilde\gamma=\partial\widetilde\Delta\)
bounds a disk containing exactly one point of \(B\), so it is
peripheral.
\end{proof}

\section{Lifting the square of the twist}

Let \(N\) be a small closed annular neighborhood of \(\gamma\).
Since \(\gamma\) has one preimage component covering it with degree
\(2\), \(f^{-1}(N)\) is an annular neighborhood of
\(\widetilde\gamma\), and
\[
f:f^{-1}(N)\longrightarrow N
\]
is a degree-\(2\) covering of annuli.

\begin{lemma}\label{lem:twist-lift}
The square \(T_\gamma^2\) admits a lift under \(f\) which is
isotopic relative to \(B\) to the identity.  More precisely,
\[
f\circ T_{\widetilde\gamma}
   =T_\gamma^2\circ f
\]
up to isotopy relative to the respective marked sets, and
\(T_{\widetilde\gamma}\) is isotopic to the identity relative to
\(B\).
\end{lemma}

\begin{proof}
Use the standard annulus model
\[
N=\{(\theta,t):\theta\in\mathbb R/2\pi\mathbb Z,\ 0\le t\le1\},
\]
and let the double cover be \((\vartheta,t)\mapsto(2\vartheta,t)\).
A twist on \(N\) is represented by
\[
(\theta,t)\longmapsto(\theta+2\pi t,t).
\]
Its square is
\[
(\theta,t)\longmapsto(\theta+4\pi t,t),
\]
which lifts to
\[
(\vartheta,t)\longmapsto(\vartheta+2\pi t,t).
\]
The latter is one full twist about \(\widetilde\gamma\).  Hence a
lift of \(T_\gamma^2\) is \(T_{\widetilde\gamma}\).

By Lemma \ref{lem:lift-peripheral}, \(\widetilde\gamma\) bounds a
disk containing only the marked point \(p_1\).  Therefore
\(T_{\widetilde\gamma}\) is isotopic to the identity relative to
\(B\).
\end{proof}

\begin{lemma}\label{lem:forget-peripheral}
After forgetting \(v\), the curve \(\gamma\) becomes peripheral
around \(p_2\).  Consequently the mapping class induced by
\(T_\gamma^2\) on \(\Teich(A)\) is trivial.
\end{lemma}

\begin{proof}
In the sphere marked only by \(A\), the point \(v\) is no longer
marked.  The disk \(\Delta\) may then be isotoped across \(v\)
without crossing any point of \(A\).  Thus \(\gamma\) is isotopic
relative to \(A\) to a small loop around \(p_2\).  A Dehn twist
about such a peripheral curve is trivial in \(\Mod(A)\).
\end{proof}

\section{Preservation of the deformation space}

We use the following standard naturality property of the Teichm\"uller
pullback.  If \(h\) is a homeomorphism of \((\Sph,B)\) and
\(\widetilde h\) is a lift under \(f\), so that
\[
f\circ\widetilde h=h\circ f,
\]
then the induced actions on Teichm\"uller space satisfy
\[
f^*(h\cdot\tau)=\widetilde h\cdot f^*(\tau),
\]
up to the conventional choice of left or right mapping-class action.
Since all twists below act trivially at the relevant marked level,
the convention is immaterial.

\begin{proposition}\label{prop:preservation}
For every \(\tau\in\D_f\),
\[
f^*(T_\gamma^2\cdot\tau)
=i_*(T_\gamma^2\cdot\tau).
\]
Hence \(T_\gamma^2(\D_f)=\D_f\).
\end{proposition}

\begin{proof}
By Lemma \ref{lem:twist-lift}, a lift of \(T_\gamma^2\) under \(f\)
is isotopic relative to \(B\) to the identity.  Therefore it acts
trivially on Teichm\"uller space, and naturality gives
\[
f^*(T_\gamma^2\cdot\tau)=f^*(\tau).
\]

On the other hand, Lemma \ref{lem:forget-peripheral} says that the
class induced by \(T_\gamma^2\) after forgetting \(v\) is trivial in
\(\Mod(A)\).  Since the forgetful map is equivariant,
\[
i_*(T_\gamma^2\cdot\tau)=i_*(\tau).
\]

Because \(\tau\in\D_f\), we have \(f^*(\tau)=i_*(\tau)\), and hence
\[
f^*(T_\gamma^2\cdot\tau)=i_*(T_\gamma^2\cdot\tau).
\]
Thus \(T_\gamma^2(\D_f)\subset\D_f\).  Applying the same argument
to \(T_\gamma^{-2}\), whose lift is a negative peripheral twist,
gives the reverse inclusion.
\end{proof}

\section{Component separation}

We now use the standard covering description of deformation spaces
and the rigidity of marked self-conjugacies.  The relevant general
results are the path-lifting theorem for deformation spaces and the
McMullen--Sullivan description of deck transformations by
quasiconformal self-conjugacies; see
\cite{McMullenSullivan}.  The relevant absence of invariant line
fields is in \cite[Chapter 3]{McMullenBook}.

\begin{lemma}[Monodromy criterion]\label{lem:monodromy}
Let \(C\) be a component of \(\D_f\), and let \(h\) be a pure mapping
class of \((\Sph,B)\) such that \(h(C)=C\).  Then \(h\) is represented
by a quasiconformal self-conjugacy of a rational map in the moduli
image of \(\D_f\), fixing \(B\) pointwise.
\end{lemma}

\begin{proof}
Choose \(\tau_0\in C\).  Since \(h(C)=C\), there is a path
\(\tau_t\) in \(C\) from \(\tau_0\) to \(h\cdot\tau_0\).
Projecting to moduli space gives a loop of conjugacy classes of
quadratic rational maps in the prescribed stratum.

By the deformation-space path-lifting theorem, this path is
represented by a continuous family of rational maps \(R_t\) and
quasiconformal conjugacies
\[
h_t:f\longrightarrow R_t
\]
preserving the labeled points of \(B\).  Normalize the family by
fixing three points of the \(c_1\)-cycle.  Since the projected path
is closed, the normalized endpoint maps satisfy \(R_1=R_0=:R\).

The endpoint markings differ by \(h\).  Hence, after changing the
endpoint representative by the corresponding isotopy, the
composition
\[
k=h_0^{-1}\circ h_1
\]
satisfies
\[
k\circ R=R\circ k
\]
and fixes the marked set \(B\).  Its isotopy class relative to \(B\)
is the class of \(h\) or of \(h^{-1}\), depending on the chosen
action convention.
\end{proof}

\begin{lemma}[Rigid marked centralizer]\label{lem:rigid-centralizer}
Let \(R\) be a quadratic rational map satisfying the hypotheses in
the statement of Theorem \ref{thm:main}.  If \(k\) is a
quasiconformal homeomorphism satisfying
\[
k\circ R=R\circ k,
\qquad k|_B=\id_B,
\]
then \(k\) is isotopic to the identity relative to \(B\).
\end{lemma}

\begin{proof}[Proof outline and citation]
This is the marked version of the standard rigidity theorem for
quasiconformal centralizers; see the deformation-theoretic
description of self-conjugacies in \cite{McMullenSullivan}.  We
indicate the dynamical steps.

Let \(\mu_k\) be the Beltrami coefficient of \(k\).  Since \(k\)
commutes with \(R\), the measurable conformal structure determined
by \(\mu_k\) is \(R\)-invariant.  McMullen's invariant-line-field
theorem \cite[Chapter 3]{McMullenBook} implies that an invariant
Beltrami field on the Julia set is zero except in the flexible
Latt\'e case.  A Latt\'e map has no periodic critical point, so the
present map is not Latt\'e; therefore \(\mu_k=0\) on the Julia set.

The remaining invariant Beltrami field is supported on the Fatou
set.  The McMullen--Sullivan straightening theorem on Fatou
components says that such a self-conjugacy is isotopic to a
conformal self-conjugacy once the marked set fixes a periodic
critical cycle and an additional critical value not belonging to
that cycle.  The latter condition removes the remaining rotations
and translations in the Fatou coordinates.  In the exceptional
power-map case, both critical values belong to one cycle of length
at most \(2\), which is incompatible with the exact period-\(4\)
cycle of \(c_1\) and the assumption on \(v\).

Thus \(k\) is isotopic relative to \(B\) to a M\"obius
transformation \(M\).  But \(M\) fixes the four distinct points of
\(A\), and a M\"obius transformation fixing three points is the
identity.  Hence \(k\) is isotopic to the identity relative to
\(B\).
\end{proof}

\begin{proposition}\label{prop:no-component-stabilized}
No nonzero power of \(T_\gamma^2\) stabilizes a component of
\(\D_f\).
\end{proposition}

\begin{proof}
Suppose that for some \(m\neq0\),
\[
T_\gamma^{2m}(C)=C.
\]
By Lemma \ref{lem:monodromy}, the mapping class \(T_\gamma^{2m}\)
is represented by a quasiconformal self-conjugacy fixing \(B\).
By Lemma \ref{lem:rigid-centralizer}, this mapping class is trivial.

On the other hand, \(\gamma\) is essential in the sphere with five
marked points: one complementary disk contains \(\{v,p_2\}\), and
the other contains \(\{p_0,p_1,p_3\}\).  A Dehn twist about an
essential simple closed curve has infinite order in the mapping
class group; see \cite[Chapter 3]{FarbMargalit}.  Consequently
\(T_\gamma^{2m}\) is nontrivial for every \(m\neq0\), a
contradiction.
\end{proof}

\begin{proof}[Proof of Theorem \ref{thm:main}]
Proposition \ref{prop:preservation} shows that \(T_\gamma^2\)
preserves \(\D_f\).  Therefore it permutes its connected components.
If
\[
T_\gamma^{2m}(C)=T_\gamma^{2n}(C),
\]
then
\[
T_\gamma^{2(m-n)}(C)=C.
\]
Proposition \ref{prop:no-component-stabilized} forces \(m=n\).
Thus all the components \(T_\gamma^{2n}(C)\), \(n\in\mathbb Z\),
are distinct.
\end{proof}

\begin{remark}
The same argument works with \(p_2\) replaced by either \(p_0\) or
\(p_3\): the enclosed disk then contains \(v\) and one regular value
of \(f\).  We chose \(p_2=f^{\circ2}(c_1)\) for concreteness.  The
essential feature is that the inverse disk contains exactly one
marked point, so the lifted square twist is peripheral, while the
downstairs twist becomes peripheral after forgetting \(v\).
\end{remark}

\begin{thebibliography}{99}

\bibitem{FarbMargalit}
B. Farb and D. Margalit,
\emph{A Primer on Mapping Class Groups},
Princeton Mathematical Series, vol.~49,
Princeton University Press, Princeton, NJ, 2012.

\bibitem{McMullenBook}
C. T. McMullen,
\emph{Complex Dynamics and Renormalization},
Annals of Mathematics Studies, vol.~135,
Princeton University Press, Princeton, NJ, 1994.

\bibitem{McMullenSullivan}
C. T. McMullen and D. P. Sullivan,
Quasiconformal homeomorphisms and dynamics III:
The Teichm\"uller space of a holomorphic dynamical system,
\emph{Advances in Mathematics} \textbf{135} (1998), no.~2,
351--395.

\end{thebibliography}

\end{document}
```